\documentclass[11pt,a4paper]{article}
\usepackage{harmonikabau_preamble}

% == Document Metadata ====================
\newcommand{\hbdoknr}{0012}
\newcommand{\hbtitel}{Reed Stiffness}
\newcommand{\hbuntertitel}{Calculation, Material Selection and Effect on Frequency}
\newcommand{\hbheadtext}{Doc.\,0012 --- Reed Stiffness}
\newcommand{\hbfoottext}{Cross-references: Doc.\,0010}

\AtBeginDocument{%
  \fancyhead[L]{\small\hbheadtext}%
  \fancyhead[R]{\small\thepage}%
  \fancyfoot[C]{\small\color{hb@darkgray}\hbfoottext}%
}

\begin{document}
\hbmaketitle

{\small\itshape The stiffness of the reed determines two properties that are in conflict: response and tuning stability. A soft reed responds more easily but detunes more strongly with bellows pressure. A stiff reed holds the pitch more steadily but requires more pressure to start vibrating.}

\section{Stiffness and Frequency}

Frequency and stiffness both depend on geometry, but to \textbf{different degrees}:
\begin{equation}
f \propto \frac{t}{L^2}, \qquad k = \frac{E \cdot W \cdot t^3}{4 \cdot L^3}
\end{equation}
Stiffness $k$ depends \textbf{cubically} on $t$, while frequency depends only linearly. 10\,\% thinner = 27\,\% softer, but only 10\,\% flatter.

\section{Detuning by Bellows Pressure}

The Bernoulli effect creates a ``negative spring'': $k_\text{eff} = k_\text{reed} - k_\text{Bernoulli}(p)$. The frequency shift: $\Delta f/f \approx -k_\text{Bernoulli}/(2 \cdot k_\text{reed})$.

\begin{keybox}
\textbf{Result:} Bass (k $\approx$ 80--150\,N/m): 4--11\,cents/kPa (audible). Treble (k $\approx$ 250--400\,N/m): 0.3--0.8\,cents/kPa (barely audible).
\end{keybox}

\section{The Trade-off}

$\Delta f \times p_\text{threshold} \approx \text{const}$ --- good response and stable tuning are physically incompatible.

\section{Tuning Weight vs.\ Thinning}

\textbf{Adding mass:} Lowers $f$, $k$ unchanged $\to$ same detuning. \textbf{Thinning:} Lowers $f$ and $k$ ($k \propto t^3$) $\to$ 37\,\% more detuning per 10\,\% thinning.

\section{Summary}
\begin{keybox}
1.~$k \propto t^3$: cubic coupling. 2.~Detuning $\propto 1/k$. 3.~$\Delta f \times p_\text{th} \approx$ const. 4.~Tuning weight is more pitch-stable than thinning. 5.~Profiling decouples mass and stiffness.
\end{keybox}

\begin{center}
\textcolor{accentred}{\rule{0.6\textwidth}{1pt}}\\[4pt]
{\small\itshape The reed seeks its equilibrium between freedom and control --- just like the player.}
\end{center}

\end{document}
